\documentclass[12pt]{article}
\usepackage[utf8]{inputenc}
\usepackage[russian]{babel}
\usepackage{amsmath}
\usepackage{amssymb}
\usepackage{geometry}
\usepackage{hyperref}
\geometry{a4paper, margin=2cm}

\begin{document}

\tableofcontents

\section{Задача 3.1.6d}

\textbf{Условие:} Отсортировать последовательность $a_1,\ldots,a_n$ в порядке возрастания элементов. Выбрана сортировка HeapSort.
\\
\texttt{in:} $n,a_1,\ldots,a_n$\\
\texttt{out:} $b_1,\ldots,b_n$, где $b_1\leq b_2\leq\ldots\leq b_n$

\section{Псевдокод}

Псевдокод HeapSort, приведённый в исходном файле:
\begin{verbatim}
public static void sort(Comparable[] pq) {
    int n = pq.length;
    for (int k = n/2; k >= 1; k--)
        sink(pq, k, n);
    int k = n;
    while (k > 1) {
        exch(pq, 1, k--);
        sink(pq, 1, k);
    }
}

private static void sink(Comparable[] pq, int l, int m) {
    while (2*l <= m) {
        int j = 2*l;
        if (j < m && less(pq, j, j+1)) j++;
        if (!less(pq, l, j)) break;
        exch(pq, l, j);
        l = j;
    }
}
\end{verbatim}

\section{Решение}

\begin{verbatim}
; Data Model
;0: r0
;1: trash | iterator
;2: &end
;3: n
;4: k
;5: l
;6: m
;7: j
;8: call (from what place was sink called last time)
;9 ... 9+n-1: arr

Load =0
Read 3
Sub 3           ; r3 := n
JZero END
JGTZ END		; n <= 0 ==> exit (invalid size)

Load 3
Add =9
Store 2			; r2 := 9 + n (arr_a.end())

Load =9
Store 1			; r1 := 9 (arr_a.begin())

READ_ARRAY: Load 2
	Sub 1
	JZero SORT		; if r2 == r1 exit read
	Read *1
	Load 1
	Add =1
	Store 1		; r1++
	Jump READ_ARRAY
SORT: Load 3
    Div =2
    Store 4     ; r4 := k = n // 2
FOR: Load =1
    Sub 4
    JGTZ FOR_END    ; if k < 1 then exit for
    Load 4
    Store 5
    Load 3
    Store 6     ; l := k, m := n
    Load =0
    Store 8     ; call = 0 (call from for)
    Jump SINK   ; call sink(k, n)
FOR_CALL_END: Load 4
    Sub =1
    Store 4     ; k--
    Jump FOR
FOR_END: Load 3
    Store 4     ; k := n
WHILE_k_gt_1: Load =1
    Sub 4
    JZero PRINT_INIT
    JGTZ PRINT_INIT   ; if k <= 1 then exit while and exit sort
    ; SWAP
    Load 4
    Add =8
    Load *0
    Store 1     ; r1 := arr[k]
    Load 9      ; r0 := arr[1]
    Store 2     ; r2 := arr[1]
    Load 1
    Store 9     ; arr[1] := arr[k]
    Load 4
    Add =8
    Store 1     ; r1 := &arr[k]
    Load 2
    Store *1    ; arr[k] (= *r1) := arr[1] (= r2)
    ; k--
    Load 4
    Sub =1
    Store 4
    ; CALL sink(1, k)
    Load =1
    Store 5     ; l := 1
    Load 4
    Store 6     ; m := k
    Load =1
    Store 8     ; call := 1 (call sink from while)
    Jump SINK
SINK: Load =2
    Mult 5
    Sub 6
    JGTZ CALL_RETURN    ; if 2*l-m > 0 then exit SINK
    Load =2
    Mult 5
    Store 7     ; j := 2*l
    ; if (j < m && less(pq, j, j+1)) j++;
    Load 7
    Sub 6
    JZero IF_BREAK
    JGTZ IF_BREAK       ; if j-m>=0 then do not do IF
    Load 7
    Add =8
    Store 1     ; r1 := &arr[j]
    Add =1
    Store 2     ; r2 := &arr[j+1]
    Load *1
    Sub *2
    JZero IF_BREAK
    JGTZ IF_BREAK       ; if a[j] >= a[j+1] then do not do IF
    Load 7
    Add =1
    Store 7     ; j++
; if (!less(pq, l, j)) break;
IF_BREAK: Load 7
    Add =8
    Load *0
    Store 1     ; r1 := arr[j]
    Load 5
    Add =8
    Load *0     ; r0 := arr[l]
    Sub 1       ; r0 := arr[l] - arr[j]
    JZero CALL_RETURN
    JGTZ CALL_RETURN    ; if a[l] >= a[j] then return from SINK
    ; SWAP if not break executed
    Load 7
    Add =8
    Load *0
    Store 1     ; r1 := arr[j]
    Load 7
    Add =8      ; r0 := &arr[j]
    Store 2     ; r2 := &arr[j]
    Load 5
    Add =8
    Load *0     ; r0 := arr[l]
    Store *2    ; arr[j] (= *r2) := arr[l] (= r0)
    Load 5
    Add =8
    Store 2     ; r2 := &arr[l]
    Load 1      ; r0 := arr[j] (previous)
    Store *2    ; arr[l] (= *r2) := arr_previous[j] (= r0)
    ; l = j
    Load 7
    Store 5
    ; goto WHILE
    Jump SINK
CALL_RETURN: Load 8
    JZero FOR_CALL_END  ; return from sink(k, n)
    Jump WHILE_k_gt_1   ; return from sink(1, k)
PRINT_INIT: Load =9
    Store 1     ; r1 := arr.begin()
    Add 3
    Store 2     ; r2 := arr.end()
PRINT: Load 2
    Sub 1
    JZero END   ; if r1 == r2 then exit print
    Write *1
    Load 1
    Add =1
    Store 1     ; r1++
    Jump PRINT
END: Halt 0
\end{verbatim}

\section{Сложность}

Пусть $h:=\lfloor n/2\rfloor$.

Всё вне вызовов SINK:
\begin{enumerate}
    \item Проверка входа - $5$ шагов
    \item Инициализация указателей - $5$ шагов
    \item Чтение массива - $8n+3$ шага
    \item SORT до FOR - $3$ шага
    \item Управление циклом построения кучи, без вызовов $\texttt{SINK}$ - $14h+5$ шагов.
    \item Обмен элементов и управление циклом извлечения максимума, без вызовов $\texttt{SINK}$ - $27n-24$ шага.
    \item Инициализация вывода, вывод массива и завершение программы - $8n+8$ шагов.
\end{enumerate}

Итого фиксированная часть:
\[
    T_{\mathrm{fixed}}(n)=43n+14\left\lfloor\frac n2\right\rfloor+5.
\]

Кроме того, выполняются $h$ вызовов $\texttt{SINK}$ при построении кучи и $n-1$ вызовов при извлечении максимума. Обозначим через $C_r(A,l,m)$ точное количество команд одного вызова $\texttt{SINK}$ для текущего состояния массива $A$, где $r=0$ означает возврат в цикл построения кучи, а $r=1$ - возврат в цикл извлечения максимума.

Точное число шагов одного вызова можно задать рекуррентно. Если $2l>m$, то
\[
    C_r(A,l,m)=6+r.
\]
Иначе положим $j=2l$. Выбранный потомок $q$ и затраты $B$ на проверку наличия правого потомка и выбор большего потомка определяются так:
\[
(q,B)=
\begin{cases}
(j,3), & j=m,\\
(j,4+8), & j<m,\ A_j=A_{j+1},\\
(j,4+9), & j<m,\ A_j>A_{j+1},\\
(j+1,4+12), & j<m,\ A_j<A_{j+1}.
\end{cases}
\]
Здесь затраты $8$, $9$ и $12$ учитывают разные исходы сравнения двух потомков; случай $12$ включает увеличение индекса потомка. Пусть $A'$ - массив после обмена элементов в позициях $l$ и $q$. Тогда
\[
C_r(A,l,m)=
\begin{cases}
7+B+p+2+r, & A_l\geq A_q,\\
7+B+10+19+C_r(A',q,m), & A_l<A_q,
\end{cases}
\]
где $p=9$, если $A_l=A_q$, и $p=10$ в остальных случаях. Эта запись считает шаги по выполненным ветвям, а не перечисляет команды по одной.

Точное число шагов для конкретного входа имеет вид
\[
\begin{aligned}
T(n,A)={}&43n+14\left\lfloor\frac n2\right\rfloor+5\\
&+\sum_{k=1}^{h} C_0(A^{(k)},k,n)
+\sum_{m=1}^{n-1} C_1(A^{[m]},1,m),
\end{aligned}
\]
где $A^{(k)}$ и $A^{[m]}$ - состояния массива непосредственно перед соответствующими вызовами $\texttt{SINK}$. Следовательно, точное число шагов в худшем случае задаётся формулой
\[
    T_{\max}(n)=\max_{A} T(n,A),
\]
где максимум берётся по допустимым входным последовательностям длины $n$.

Количество шагов внутри $\texttt{SINK}$ зависит от данных: от того, сколько уровней проходит элемент, имеется ли правый потомок и какой из двух потомков выбирается. Поэтому для вычисления максимума требуется отдельно максимизировать эти ветвления с учётом изменений массива между вызовами. Ниже приведена явная верхняя оценка, достаточная для определения порядка роста.

Для оценки худшего случая заметим следующее:
\begin{itemize}
    \item При построении кучи число переходов вниз не превышает
    \[
        H(n)=\sum_{k=1}^{\lfloor n/2\rfloor}\left\lfloor\log_2\frac nk\right\rfloor
        =n-s_2(n),
    \]
    где $s_2(n)$ - число единиц в двоичной записи $n$.
    \item При извлечении максимума число переходов вниз не превышает
    \[
        D(n)=\sum_{m=1}^{n-1}\lfloor\log_2 m\rfloor.
    \]
    \item Один переход вниз вместе с соответствующими проверками занимает не более $52$ шагов, а завершающая часть одного вызова $\texttt{SINK}$ - не более $36$ шагов.
\end{itemize}

Следовательно, получаем явную верхнюю оценку:
\[
\begin{aligned}
T_{\max}(n)\leq{}&43n+14\left\lfloor\frac n2\right\rfloor+5\\
&+52\bigl(H(n)+D(n)\bigr)
+36\left(\left\lfloor\frac n2\right\rfloor+n-1\right).
\end{aligned}
\]
Так как $H(n)=O(n)$ и $D(n)=O(n\log n)$, худшее время работы HeapSort составляет $O(n\log n)$. Для сортировки сравнением в худшем случае необходимо $\Omega(n\log n)$ сравнений, поэтому
\[
    \boxed{T_{\max}(n)=\Theta(n\log n)}.
\]
При стандартном предположении о равномерно случайной перестановке различных элементов среднее время HeapSort также составляет $\Theta(n\log n)$.

\end{document}